% =============================================================================
% Template: Trabalho Acadêmico — UnB Modelo A (engenharia)
%
% Capa com logo grande e faixa azul institucional.
%
% Compilação local (a partir da raiz do repositório):
%   TEXINPUTS=./src: latexmk -pdf templates/unb-A/main.tex
% =============================================================================

\documentclass{abnt-unb-A}

\usepackage{abnt-resumo}
\usepackage{abnt-citacoes}   % carrega abnt-refs automaticamente

\addbibresource{refs.bib}

% --- Metadados (preencha com seus dados) ---
\titulo{Título do Trabalho}
\subtitulo{subtítulo, se houver}
\autor{Nome Completo do Autor}
\orientador{Prof. Dr. Nome do Orientador}
% \coorientador{Prof. Dr. Nome do Coorientador}
\faculdade{Faculdade de Tecnologia}
\departamento{Departamento de Engenharia Elétrica}
\curso{Engenharia Elétrica}
\grau{Bacharel em Engenharia Elétrica}
\tipotrabalho{Trabalho de Conclusão de Curso}
\unblogo{unb_logo}
\ano{2026}
\natureza{Trabalho de Conclusão de Curso apresentado ao Departamento
  de Engenharia Elétrica da Universidade de Brasília como requisito
  parcial para obtenção do grau de Bacharel em Engenharia Elétrica.}
\programa{Graduação em Engenharia Elétrica}

% =============================================================================
\begin{document}

% --- Pré-textuais ---
\imprimircapa
\imprimirfolhaderosto

\dataaprovacao{1 de janeiro de 2026}
\membrodabanca{Prof. Dr. Membro Um}{ENE/UnB}
\membrodabanca{Prof. Dr. Membro Dois}{ENE/UnB}
\membrodabanca{Prof. Dr. Membro Três}{ENE/UnB}
\imprimirfolhadeaprovacao

\imprimirdedicatoria{%
  Dedicatória (opcional).}

\imprimiragradecimentos{%
  Agradecimentos (opcional).}

\imprimirresumo{%
  Resumo em português.%
}{Palavra-chave 1. Palavra-chave 2. Palavra-chave 3}

\imprimirabstract{%
  Abstract in English.%
}{Keyword 1. Keyword 2. Keyword 3}

\imprimirsumario

% --- Textuais ---
\chapter{Introdução}

Texto da introdução \cite{eco1977}.

\chapter{Desenvolvimento}

Texto do desenvolvimento \cite{medeiros2019}.

\chapter{Conclusão}

Texto da conclusão.

% --- Pós-textuais ---
\postextual
\printbibliography

\end{document}
